\section{An Introduction to Haskellog}

\subsection{Goals}

Haskellog is an interpreter for a logic programming language similar to Prolog (with a few functionalities missing and a slightly different syntax). The user gives a semantic meaning to some functions, and the program can then answer queries: it can determine whether a certain formula can be derived from the specified rules and, if it can, provide an instantiation of the variables that makes it derivable.

The program expects a file containing logical rules made from terms.
Terms are made with functions in lower case, variables in upper case, constants, which are simply functions of arity 0 and integers which are just syntactic sugar for Peano arithmetic: $0 := z(),\ 1 := s(0),\ 2 := s(s(0)),\ \dots$
Each rule is either a fact, which is a single term like \verb|plus(X,z,X).| that you could interpret as $\forall X \in \mathbb{N}, X+0 = X$ or a rule composed of a conclusion and some premises, such as \verb|sum_of_consecutives(N) :- plus(A, 1, B), plus(A,B,N).| which could be interpreted as "a number $N$ is a sum of two consecutive numbers $A$ and $B$ if $A + B = N$ and $A + 1 = B$". Obviously these interpretations are only given by the rules: the program does not know what \verb|plus| means to you and if you define it, it is your responsability to provide rules that match your intended meaning of each predicate.
\newline
For example, we can enrich the semantics of \verb|plus| that we started defining before so that it matches the idea that \verb|plus(X,Y,Z)| should be derivable if $X+Y=Z$ by adding the rule \verb|plus(X,s(Y),s(Z)) :- plus(X,Y,Z).|.

The program can then answer queries, which are a succession of terms interpreted as a conjunction. Continuing with the previous example, a query could be \verb|?- plus(X,Y,4), plus(X,2,Y)| which receives a positive answer from the program: \verb|yes, X = 1 and Y = 3|. Once Haskellog has output one possible answer, the user can ask for additional ones.\newline You might think that the query \verb|?- plus(X,Y,4), plus(X,1,Y)| would be given a similar answer \verb|yes, X = 1,5 and Y = 2,5| but remember that we never defined what 1,5 or 2,5 are, the semantics only talks about integers. Therefore, the answer would be \verb|no|. Try implementing half-integers to change that!

Internally, Haskellog works as follows: a parser reads the source file and transforms the rules and queries in some Haskell data types. Then, when a query is inputed in the UI, it is parsed and the prover attempts to solve it by resolution: the query is viewed as a conjunction of goals, and Haskellog repeatedly selects one goal and tries to match it with the head of a fact or rule from the input file. This matching is performed by unification, which computes a substitution for variables making two terms identical when possible. If no such substitution exists, that branch of the search fails and the program backtrack. When a rule matches, its premises replace the resolved goal, and the process continues recursively until either all goals have been solved, yielding a successful proof and an answer substitution, or no rule applies. The result of the proof is then presented to the user.

\subsection{Usecase: Proof Theory Homework Checker}
One could use Haskellog to see if rule (in a proof theoretic sense) is deliverable. 
